\documentclass[10pt,a4paper]{amsart}
\usepackage[margin=19mm]{geometry}
\usepackage{amsmath,amssymb,mathtools}
\usepackage[scr=rsfs,cal=euler]{mathalfa}
\usepackage{xfrac}
\usepackage[unicode,hidelinks,bookmarks=false]{hyperref}
\newcommand{\ie}{i.e.\ }
\newcommand{\cf}{cf.\ }
\newcommand{\resp}{respectively\ }
\newcommand{\Subsection}{\subsection}
\providecommand{\nobreakdash}{\nobreak\hbox{-}}
\setlength{\emergencystretch}{2em}
\multlinegap0pt
\usepackage{mdframed}
\usepackage{mdframed}
\usepackage{mdframed}
\usepackage{mdframed}
\usepackage{amssymb}
\usepackage{mathtools}
\usepackage{tikz}
\usepackage{mdframed}
\usepackage{tcolorbox}
\usepackage{xcolor}
\usepackage[shortlabels]{enumitem}
\newtheorem{proposition}{Proposition}
\newtheorem{corollary}{Corollary}
\newtheorem{lemma}{Lemma}
\newtheorem{claim}{Claim}
\newtheorem{proofclaim}{Proofclaim}
\usepackage{cleveref}
\crefname{proposition}{Proposition}{Propositions}
\crefname{corollary}{Corollary}{Corollarys}
\crefname{lemma}{Lemma}{Lemmas}
\crefname{claim}{Claim}{Claims}
\crefname{proofclaim}{Proofclaim}{Proofclaims}
\newcommand{\Bin}{\operatorname{Bin}}
\renewcommand{\Pr}{\mathbb{P}}
\newcommand{\SMA}{\operatorname{SMALL}}
\def\calP{\mathcal{P}}
\newcommand{\eps}{\varepsilon}
\title{On the hitting time of Hamiltonicity in bipartite Dirac graphs}
\author{Yiting Wang}
\date{Source: arXiv:2606.11992v1 (CC BY 4.0)}
\begin{document}
\maketitle

\noindent\emph{Source and licence.} On the hitting time of Hamiltonicity in bipartite Dirac graphs. Version arXiv:2606.11992v1 (CC BY 4.0), distributed by arXiv under the Creative Commons Attribution 4.0 International licence, http://creativecommons.org/licenses/by/4.0/, as recorded in the arXiv licence metadata for this exact version and shown on the abstract page https://arxiv.org/abs/2606.11992v1. Full licence notice: \texttt{licenses/arxiv-2606.11992-CC-BY-4.0.txt}.\par\smallskip

\noindent\textit{Editorial scope.} Counted unit: the article's lemma lem:sparse-expander-exists (source0.tex lines 494--496). The article's complete original proof of it is delivered with it (source0.tex lines 497--585, one \texttt{proof} environment of 89 lines). No mathematics is written, repaired or re-derived: the statement and the proof are the article's own lines, in the article's own notation, and no retained line is substituted or re-flowed.  Retained uncounted as the source-local context the proof needs: lines 320--322; lines 407--409; lines 417--425.  Nothing was omitted from any retained span, so the package carries no omission record.  No retained line contains a citation, so the wrapper carries no bibliography and no bibliography entry.  No retained line contains a figure environment, an image-inclusion command or drawing code, so no image is delivered and the package declares no asset dependency.  Editorial formatting only, in addition: an AMS \texttt{amsart} preamble that restates the article's own declarations and macros used by the retained lines, namely the environments \texttt{proposition}, \texttt{corollary}, \texttt{lemma}, \texttt{claim}, \texttt{proofclaim} on separate counters, together with the standard packages those lines need.  The retained environments print in this document as Proposition 1, Corollary 1, Lemma 1, Claim 1, Proofclaim 1 in order of appearance, whereas the article prints the same items with the numbering of its own full document.  The complete native source of the article as distributed in the arXiv e-print for this exact version is bundled unchanged as \texttt{sources/202594/source0.tex}.
\begin{proposition}\label{prop:basic-inequality}
    Let $X\sim \Bin(n,p)$. For any $t > 0$, we have $\Pr(X\geq t)\leq \binom{n}{t}p^t$. For any $t< np$, we have $\Pr(X\leq t)\leq t\binom{n}{t}p^t(1-p)^{n-t}$.
\end{proposition}

\begin{corollary}\label{cor:monotonicity}
    Let $\calP$ be a monotone decreasing property. Then, if $G^+\in \calP$ whp, then so is $G_\tau$. Similarly, for a monotone increasing property $\calP$, if $G^-\in \calP$ whp, then so is $G_\tau$.
\end{corollary}

\begin{lemma}\label{lem:small-property}
        Whp, the following properties hold:
        \begin{enumerate}
            \item $\Delta(G_\tau)\leq 20\log n$,
            \item $|\SMA(G_\tau)|\leq n^{(1-\eps)/2}$,
            \item For every $u\in \SMA(G_\tau)$, there is no cycle of length at most $4$ containing $u$ in $G_\tau$,
            \item For distinct vertices $u,v\in \SMA(G_\tau)$, there is no path of length at most $4$ between them in $G_\tau$.
        \end{enumerate}
    \end{lemma}

\begin{lemma}\label{lem:sparse-expander-exists}
    Whp $G_\tau$ contains a connected spanning subgraph $R$ such that $R$ is an $(\alpha n,2)$-expander and $e(R)\leq 3 \xi n \log n$.
\end{lemma}

\begin{proof}
    Recall that $\SMA(H) = \{v\in V(H): d_{H}(v)\leq \xi \log n\}$. 
    Let $G_1\subseteq G_\tau$ be obtained as follows: For each vertex $v\in V(G)\setminus \SMA(G_\tau)$, choose a uniformly random set of $\xi\log n$ incident edges. 
    We first analyse the expansion properties of $G_1$:
    \begin{claim}\label{claim:4-expander}
        Whp, for every $S\subseteq V(G)\setminus \SMA(G_\tau)$ such that $S\subseteq A$ or $S\subseteq B$ and $|S|\leq \alpha n$, we have $|N_{G_1}(S)|\geq 4|S|$.
    \end{claim}
    \begin{proofclaim}
        Suppose for contradiction that there exists a set $S\subseteq A\setminus \SMA(G_\tau)$ of size at most $\alpha n$ such that its neighbourhood $T = N_{G_1}(S)\subseteq B$ has size at most $4|S|$ (the symmetric case where $S\subseteq B$ and $T\subseteq A$ can be handled similarly).
        We distinguish cases according to $|S|$. 
        \begin{itemize}
            \item If $|S|\leq \xi \log n/4$, note that for any $v\in S$, we have $|T|\geq d_{G_1}(v) \geq \xi \log n\geq 4|S|$. 
            \item If $\xi\log n/4\leq |S|\leq n/\sqrt{\log n}$, note that $e(G_\tau[S\cup T]) \geq e(G_1[S\cup T])\geq \xi|S|\log n$ and $|S\cup T|\leq 5|S|$. 
            We now bound the probability that such a set $S\cup T$ exists in $G_\tau$.
            Using~\Cref{cor:monotonicity}, we do the computation in $G^+$.
            By~\Cref{prop:basic-inequality} and the union bound, the probability that there exists such a set $S\cup T$ is at most
            \begin{multline*}
                \sum_{i= 5\xi\log n/4}^{5n/\sqrt{\log n}}\binom{n}{i}\cdot \Pr\left(\Bin\left(\binom{i}{2},p_+\right) \geq \xi i\log n/5\right)
                \leq \sum_{i= 5\xi\log n/4}^{5n/\sqrt{\log n}}\left(\frac{en}{i}\right)^i \binom{i^2/2}{\xi i\log n/5} p_+^{\xi i\log n/5} = o(1)\,.
            \end{multline*}
            \item If $n/\sqrt{\log n}\leq |S|\leq \alpha n$, then there is no edge between $S$ and $B\setminus T$ in $G_1$. 
            We first bound the probability that there exists a set $S$ of this size such that $e_{G_\tau}(S,B\setminus T)< n\sqrt{\log n}/20$.
            By~\Cref{cor:monotonicity}, we do the computation in $G^-$.
            
           Note that there are at least $(1/2 + \eps - 4\alpha)n\cdot n/\sqrt{\log n} = (1/10+\eps)n^2/\sqrt{\log n}\geq n^2/(10\sqrt{\log n})$ edges between $S$ and $B\setminus T$ in $G$.
            Each edge is included in $G^-$ independently with probability $\log n/n$. Thus, by Chernoff's bound and the union bound over all $S$ and $B\setminus T$, we have $e_{G^-}(S,B\setminus T)\geq n\sqrt{\log n}/20$ whp. 
            
            Let $F = E_{G_\tau}(S,B\setminus T)$, and for each vertex $v\in S$, let $q_v$ denote the number of edges of $F$ incident to $v$, so that $\sum_{v\in S}q_v = |F|\geq n\sqrt{\log n}/20$ whp.
            Now we bound the probability that none of $F$ appears in $G_1$. 
            Since $\Delta(G_\tau)\leq 20\log n$ whp by~\Cref{lem:small-property}, 
            we deduce that the probability that none of the $\xi\log n$ edges selected by $v$ contain an edge from $F$ is at most 
            \[
                \frac{\binom{\deg_{G_\tau}(v)-q_v}{\xi\log n}}{\binom{\deg_{G_\tau}(v)}{\xi\log n}}\leq \exp(-\xi\log n\cdot q_v/\deg_{G_\tau}(v))\leq \exp(-\xi q_v/20)\,.
            \]
            Note that the edge sets selected by distinct vertices $v\in S\subseteq A$ are disjoint, so the events concerning these vertices are independent. Thus, 
            \[
                \Pr(E_{G_1}(S,B\setminus T) = \emptyset)\leq \exp\left(-\frac{\xi}{20}\sum_{v\in S}q_v\right) = \exp(-\Omega(n\sqrt{\log n}))\,.
            \]
            The union bound over all $S$ and $T$ finishes the proof. 
        \end{itemize}
        Thus, the statement holds in all cases.
    \end{proofclaim}
    Let $G_2\subseteq G_\tau$ be the graph obtained by taking all edges incident to $\SMA(G_\tau)$ and let $H = G_1\cup G_2$. 
    We now show the expansion property of $H$:
    \begin{claim}\label{cla:G-expander}
        Whp $H$ is an $(\alpha n,2)$-expander.
    \end{claim}
    \begin{proofclaim}
        Suppose, for contradiction, that there exists a set $S\subseteq A$ of size at most $\alpha n$ such that its neighbourhood $T = N_{H}(S)\subseteq B$ has size at most $2|S|$. 
        Let $S_1 = S \cap \SMA(G_\tau)$ and $S_2 = S\setminus \SMA(G_\tau)$. 
        Note that $|N_{H}(S_1)| \geq 2|S_1|$. 
        Indeed, at time $\tau$, every $u\in S_1$ has at least two $H$-neighbours. If $|N_H(S_1)|<2|S_1|$, this would contradict~\Cref{lem:small-property}(4) by giving length-$2$ paths between two vertices of $\SMA(G_\tau)$.
        By~\Cref{claim:4-expander}, we know $|N_{H}(S_2)|\geq 4|S_2|$. By~\Cref{lem:small-property}, for any $v\in S_2$, we know that $|N_{H}(S_1)\cap N_{H}(v)|\leq 1$.
        Indeed, by~\Cref{lem:small-property}(3), for each $u \in S_1$, we have $|N_H(u)\cap N_H(v)|\leq 1$; otherwise, we obtain a $4$-cycle containing a $\SMA(G_\tau)$ vertex. Moreover, suppose there are distinct $u,u'\in S_1$ such that $|N_H(u)\cap N_H(v)| = |N_H(u')\cap N_H(v)| =  1$. Then there exists a path of length $4$ between $u$ and $u'$, contradicting~\Cref{lem:small-property}(4).
        Thus, $$|N_H(S)| = |N_H(S_1)| + |N_H(S_2)| - |N_H(S_1)\cap N_H(S_2)| \geq 2|S_1| + 4|S_2| - |S_2| \geq 2|S|\,,$$
        as desired.
    \end{proofclaim}
    From~\Cref{cla:G-expander}, it is immediate that every component of $H$ has size at least $4\alpha n$, with at least $2\alpha n$ vertices on each side.
    Thus, the total number of components of $H$ is at most $1/(2\alpha) = O(1)$.
    We now show that there is a small set of edges that ``merge'' all components of $H$:
    \begin{claim}
        Whp, there exists a set of edges $T\subseteq E(G_\tau)$ of size at most $1/(2\alpha)$ such that $H\cup T$ is connected.
    \end{claim}
    \begin{proofclaim}
        Let $\ell\leq 1/(2\alpha)$ be the number of components of $H$ and let $C_1,\dots, C_\ell$ be an enumeration of the components. Moreover, let $A_i = C_i\cap A$ and $B_i = C_i\cap B$ for each $i\in [\ell]$. Using the $(\alpha n,2)$-expansion, we know $|A_i|,|B_i|\geq 2\alpha n$ for all $i\in [\ell]$.

        Define an auxiliary component graph $G_{\mathrm{comp}}$ as follows: the vertices of $G_{\mathrm{comp}}$ are $\{C_1,\dots, C_\ell\}$, and we put an edge between components $i,j\in V(G_{\mathrm{comp}})$ if there exists an edge between $C_i$ and $C_j$ in $G_\tau$. The goal is to show that $G_{\mathrm{comp}}$ is connected, so that we can find a spanning tree in $G_{\mathrm{comp}}$. Replacing each edge of the spanning tree with an edge in $G_\tau$ that connects the corresponding two components gives the desired $T$.

        Suppose for contradiction that $G_{\mathrm{comp}}$ is not connected. Take a non-empty proper union of components and write $X\subseteq A$ and $Y\subseteq B$ for its two sides. 
        Since both the union and its complement contain whole components, we know $|X|, |A\setminus X|, |Y|, |B\setminus Y|\geq 2\alpha n$. If there is no $G_\tau$ edge crossing the component cut, then 
        \begin{equation}\label{eqn:no-edge}
            e_{G_\tau}(X, B\setminus Y) + e_{G_\tau}(A\setminus X, Y) = 0\,.
        \end{equation}
        However, using $\delta(G) \geq (1/2 + \eps)n$, we conclude 
        \[
            e_G(X,B\setminus Y) \geq |X|\cdot \max\{(1/2+\eps)n - |Y|, 0\} \text{ and } e_G(A\setminus X, Y) \geq |A\setminus X|\cdot \max\{|Y| - (1/2-\eps)n, 0\}\,.
        \]
        By a case distinction on $|Y|$, we have $e_G(X,B\setminus Y) + e_G(A\setminus X, Y)= \Omega(n^2)$. 
        We claim there is no such $X$ and $Y$ satisfying~\eqref{eqn:no-edge}. Indeed, by~\Cref{cor:monotonicity}, we may do the computation in $G^-$ and then
        \[
            \Pr(e_{G^-}(X, B\setminus Y) + e_{G^-}(A\setminus X, Y) = 0)\leq \Pr(\Bin(\Omega(n^2), \log n/n) = 0) = o(4^{-n})\,.
        \]
        Applying the union bound over all choices of $X$ and $Y$ proves that, whp, there is at least one edge across each component cut, and so $G_{\mathrm{comp}}$ is connected.
    \end{proofclaim}    
  Let $R = H\cup T$. It remains to show that $R$ is the desired graph. Note that $R$ is connected and spanning by definition and is an $(\alpha n,2)$-expander by~\Cref{cla:G-expander}.
  Also, note that 
  $$e(R)\leq e(G_1) + e(G_2) + e(T) \leq 2\xi n\log n + \xi \log n\cdot n^{(1-\eps)/2} + O(1)\leq 3\xi n\log n\,,$$
  as desired.
\end{proof}

\end{document}
